\documentclass[11pt,a4paper]{article}

% ============================================================
% PAQUETES
% ============================================================
\usepackage[utf8]{inputenc}
\usepackage[T1]{fontenc}
\usepackage[spanish,es-noshorthands]{babel}
\usepackage{amsmath,amssymb}
\usepackage{graphicx}
\usepackage{geometry}
\geometry{top=2.2cm, bottom=2.2cm, left=2.2cm, right=2.2cm}
\usepackage{float}
\usepackage{enumitem}
\setlist{nosep, leftmargin=1.5em}
\usepackage{parskip}
\setlength{\parindent}{1.25em} % Sangría de primera línea
\usepackage{indentfirst} % Sangra también el primer párrafo tras cada sección
\usepackage{caption}
\usepackage{url}
\usepackage{xcolor}
\usepackage{booktabs}
\usepackage{tabularx}
\usepackage{array}
\newcolumntype{L}{>{\raggedright\arraybackslash}X}

% Configuración de hyperref (combinando el color azul de enlaces con el estilo sin marcos)
\usepackage{hyperref}
\hypersetup{
	colorlinks=true,
	linkcolor=black,
	urlcolor=blue,
	citecolor=black,
	pdfborder={0 0 0}
}

% Renombrar cuadros a Tablas en español
\addto\captionsspanish{\renewcommand{\tablename}{Tabla}}

% Títulos de sección más compactos
\makeatletter
\renewcommand\section{\@startsection{section}{1}{\z@}%
  {-3ex \@plus -.5ex \@minus -.2ex}{1.2ex \@plus .2ex}{\normalfont\large\bfseries}}
\renewcommand\subsection{\@startsection{subsection}{2}{\z@}%
  {-2.2ex \@plus -.5ex \@minus -.2ex}{0.8ex \@plus .2ex}{\normalfont\normalsize\bfseries}}
\makeatother

% Entorno thebibliography sin el encabezado automático "Referencias":
% ya ponemos nuestro propio \section*{Bibliografía} antes de usarlo.
\makeatletter
\renewenvironment{thebibliography}[1]
 {\list{\@biblabel{\@arabic\c@enumiv}}%
       {\settowidth\labelwidth{\@biblabel{#1}}%
        \leftmargin\labelwidth
        \advance\leftmargin\labelsep
        \itemsep 1.5pt \parsep 0pt
        \usecounter{enumiv}%
        \let\p@enumiv\@empty
        \renewcommand\theenumiv{\@arabic\c@enumiv}}%
  \sloppy\clubpenalty4000\widowpenalty4000%
  \sfcode`\.\@m}
 {\def\@noitemerr
   {\@latex@warning{Empty `thebibliography' environment}}%
  \endlist}
\makeatother
\captionsetup{
    font={small,it},
    labelfont={small,it,bf},
    justification=centering,
    singlelinecheck=false
}

% ============================================================
% DATOS DEL TRABAJO
% ============================================================
\newcommand{\materia}{Procesamiento de Señales e Imágenes Biomédicas}
\newcommand{\codigoMateria}{16.63}
\newcommand{\tituloTrabajo}{Anteproyecto}
\newcommand{\subtituloTrabajo}{Clasificación automática de estadios de sueño a partir de la polisomnografía}
\newcommand{\integranteUno}{65555 -- Cassano Nieto, Facundo Tomás (\href{mailto:fcassanonieto@itba.edu.ar}{fcassanonieto@itba.edu.ar})}
\newcommand{\integranteDos}{65414 -- D'Elía, Lucas Valentín (\href{mailto:ldelia@itba.edu.ar}{ldelia@itba.edu.ar})}
\newcommand{\integranteTres}{65560 -- Ortega, Facundo (\href{mailto:facortega@itba.edu.ar}{facortega@itba.edu.ar})}
\newcommand{\comision}{BIO}
\newcommand{\docenteUno}{Tomás, Roberto Sebastián}
\newcommand{\docenteDos}{Capurro, Lucila}
\newcommand{\docenteTres}{Reppucci, Victoria}
\newcommand{\docenteCuatro}{Tassara, Francisco Joaquín}
\newcommand{\fechaEntrega}{10 de octubre de 2026}
\newcommand{\logoITBA}{logo_itba.png} % colocar el archivo del logo en esta carpeta con este nombre

\begin{document}

% ============================================================
% CARÁTULA
% ============================================================
\begin{titlepage}
    \centering
    \vspace*{1cm}
    \IfFileExists{\logoITBA}{%
        \includegraphics[width=0.45\textwidth]{\logoITBA}\par
        \vspace{0.1cm}
    }{}
    {\LARGE Instituto Tecnológico de Buenos Aires\par}
    \vspace{0.3cm}
    {\Large \materia\ (\codigoMateria)\par}
    \vspace{0.3cm}
    {\large Comisión \comision\par}
    \vfill
    {\Huge\bfseries \tituloTrabajo\par}
    \vspace{0.15cm}
    {\LARGE \subtituloTrabajo\par}
    \vfill
    {\large\bfseries Grupo 2\par}
    \vspace{0.3cm}
    {\large \integranteUno \par}
    {\large \integranteDos \par}
    {\large \integranteTres \par}
    \vfill
    {\large\bfseries Docentes\par}
    \vspace{0.3cm}
    {\large \docenteUno \par}
    {\large \docenteDos \par}
    {\large \docenteTres \par}
    {\large \docenteCuatro \par}
    \vfill
    {\large \fechaEntrega \par}
\end{titlepage}

% ============================================================
% 1. INTRODUCCIÓN
% ============================================================
\section{Introducción}

\subsection{Problemática}

La polisomnografía (PSG) es el estudio de referencia para diagnosticar
trastornos del sueño: durante una noche se registran en simultáneo EEG,
EOG, EMG de mentón y ECG. Según el
manual de la American Academy of Sleep Medicine (AASM) \cite{aasm}, el
registro se divide en épocas de 30~s y a cada una se le asigna uno de
cinco estadios: vigilia (W), sueño ligero (N1 y N2), sueño profundo (N3)
o sueño REM. La secuencia de estadios de toda la noche se denomina
hipnograma.

Hoy esta clasificación la hace un técnico a mano, época por época (unas
960 por noche), y el resultado depende de quién lo haga. El acuerdo entre dos evaluadores se
mide con el coeficiente $\kappa$ de Cohen, que vale 1 si coinciden en
todas las épocas y 0 si coinciden solo lo que se esperaría por azar.
Entre técnicos de distintos centros se reportó $\kappa = 0{,}76$
\cite{danker} y entre los técnicos del Haaglanden Medisch Centrum (HMC),
el hospital que generó la base de datos que usaremos, $\kappa = 0{,}74$
\cite{alvarez}. N1 es el estadio en el que menos coinciden.

\subsection{Marco teórico}

\textbf{Qué información aporta cada señal.} A medida que el sueño se
profundiza, el EEG pasa de una actividad rápida y de baja amplitud en
vigilia (ritmos $\alpha$, 8--13~Hz, y $\beta$, 13--30~Hz) a ondas lentas
de gran amplitud en N3 (ritmo $\delta$, 0,5--4~Hz). En N2 aparecen
además los husos de sueño, que son ráfagas de 11 a 16~Hz que duran entre 0,5 y
2~s. El sueño REM tiene un EEG parecido al de N1, por lo que se
reconoce con otras dos señales: los movimientos oculares rápidos, que se
ven en el EOG, y la pérdida de tono muscular, que se ve en el EMG del
mentón. El corazón también cambia con el estadio, ya que la variabilidad de la
frecuencia cardíaca (HRV) depende del sistema nervioso autónomo, que pasa a un predominio
parasimpático en el sueño profundo y a uno simpático en REM \cite{trinder}.

\textbf{Artefactos.} Los movimientos de los ojos producen ondas lentas y grandes en los
electrodos frontales que pueden confundirse con
sueño profundo, y la actividad del corazón aparece en el EEG como un
pulso que se repite en cada latido. Ambos ocupan las mismas frecuencias
que los ritmos de interés, por lo que un pasa-banda no puede filtrarlos
sin quitar también información útil.

\subsection{Estado del arte}

Los métodos de clasificación automática se agrupan en dos familias. La
primera calcula descriptores elegidos por el diseñador, como la potencia
del EEG en cada banda de frecuencia, y se los entrega a un
clasificador; un ejemplo es YASA \cite{yasa}, que con este enfoque logra
un acuerdo comparable al de los técnicos. La segunda usa redes
neuronales profundas que reciben la señal cruda y aprenden solas qué
patrones usar; suelen rendir algo mejor, pero no permiten saber en qué
se basa cada decisión. El artículo propuesto por la cátedra
\cite{alvarez} pertenece a esta familia, es una red convolucional y
recurrente que alcanzó $\kappa = 0{,}79$ sobre los datos de testeo.

Nuestro trabajo pertenece a la primera familia, porque permite explicar
cada decisión a partir de propiedades medibles de las señales. Además de
las técnicas vistas en clase, incorporamos tres nuevas: la estimación espectral
\emph{multitaper} \cite{thomson,prerau}, para obtener espectros menos
ruidosos de cada época; el promediado de Woody \cite{woody}, para
estimar la forma de las ondas lentas aunque no estén bien alineadas; y
los patrones espaciales comunes (CSP) \cite{ramoser}, que entra dentro del análisis
espacial y nos permite combinar los cuatro canales de EEG de la forma que mejor separe los estadios que más
se confunden.

\subsection{Objetivos e hipótesis}

\textbf{Objetivo general.} Desarrollar un algoritmo que clasifique cada
época de una PSG en uno de los cinco estadios del sueño, con decisiones
explicables a partir de las señales, que alcance un acuerdo con
los técnicos cercano al que tienen los técnicos entre sí
($\kappa \approx 0{,}74$), junto con una interfaz gráfica que muestre
cada etapa del procesamiento.

\textbf{Hipótesis.} \textbf{H1}, el espectro multitaper varía menos que el de
Welch entre épocas de un mismo estadio; \textbf{H2}, el promediado de
Woody recupera la forma de la onda lenta con menor error que el promedio
ordinario y ayuda a separar N2 de N3; \textbf{H3}, los descriptores CSP
mejoran la clasificación de N1 y REM; \textbf{H4}, quitar los artefactos
y agregar la HRV mejora el acuerdo con los técnicos.

% ============================================================
% 2. MATERIALES Y MÉTODOS
% ============================================================
\section{Materiales y métodos}

\subsection{Base de datos}

Usaremos la base de datos \textbf{HMC Sleep Staging Database v1.1} \cite{hmc}
(opción 2 de la consigna). Contiene 148 PSG de adultos derivados al centro por
distintos trastornos del sueño, registradas con equipos SOMNOmedics y
electrodos de Ag/AgCl a 256~Hz, en formato EDF.
Cada registro tiene cuatro canales de EEG (F4, C4, O2 y C3, referidos a
la mastoide del lado opuesto), dos de EOG, EMG de mentón y ECG, y está
anotado por técnicos según el manual AASM vigente. La elegimos porque
tiene los cuatro EEG que necesita CSP, el EOG que sirve de referencia
del artefacto ocular y el ECG para detectar los latidos.

Comenzaremos trabajando con 60 a 80 registros elegidos al azar para reducir el costo computacional,
y separaremos la totalidad de datos entre entrenamiento (80\,\%) y prueba (20\,\%),
obligando a que todas las épocas de un paciente quedan del mismo
lado para que no se mezclen los mismos registros en los datos de testing y de validación.

\subsection{Metodología}

\begin{enumerate}[label=\textbf{\arabic*.}, leftmargin=1.6em, itemsep=6pt, topsep=4pt]
\item \textbf{Análisis exploratorio.} Espectro de cada canal por
  estadio y verificación de artefactos, entre ellos, correlación entre EEG frontal y
  EOG para el ocular, y promedio del EEG alrededor de cada latido para
  el cardíaco (si aparece una forma repetible, hay artefacto).

\item \textbf{Filtrado convencional.} Un pasa-altos en 0,5~Hz para eliminar la
  deriva lenta de la línea de base (por sudoración o respiración) y un
  notch en 50~Hz para eliminar la interferencia de la red eléctrica de Países
  Bajos, que es el país donde se tomaron los datos; la frecuencia de corte superior se fijará a partir del espectro
  de la etapa 1. Compararemos un diseño FIR y uno IIR. El IIR necesita un
  orden mucho menor, pero desplaza de forma distinta cada frecuencia y
  deforma las ondas lentas (que se soluciona aplicandolo dos veces, una
  hacia adelante y otra hacia atrás, de modo que se compensen los desfasajes).

\item \textbf{Remoción del artefacto ocular.} Usamos el EOG como copia
  del artefacto y diseñamos un filtro adaptativo que reste la parte del
  EOG que aparece en cada canal de EEG. Se usa la diferencia entre los dos canales de EOG,
  porque en ella los movimientos oculares se refuerzan y la actividad
  cerebral, que llega igual a ambos, se cancela. Se lo comparará con el filtro de Wiener, que calcula una corrección fija.

\item \textbf{Remoción del artefacto cardíaco.} Con el algoritmo de
  Pan-Tompkins se detecta cada latido en el ECG. El
  artefacto tiene la misma forma en cada latido, mientras que el EEG no
  está sincronizado con el corazón; por eso, al promediar el EEG
  alrededor de los últimos 30 latidos el EEG se cancela y queda la forma
  del artefacto, que se resta en cada latido. Los instantes de los
  latidos se reutilizan para la HRV.

\item \textbf{Segmentación.} Épocas de 30~s, dentro de las cuales se
  supone que la señal no cambia sus propiedades. Con una segmentación
  adaptativa con un modelo autorregresivo se puede marcar las épocas
  donde eso no se cumple (por ejemplo, por un despertar breve).

\item \textbf{Espectro multitaper.} El periodograma de cada época es muy
  ruidoso. El método de Welch visto en clase lo suaviza partiendo la
  época en trozos y promediando sus espectros, pero cada trozo es más
  corto y se pierde resolución en frecuencia. El método multitaper
  usa la época completa, la multiplica por varias
  ventanas distintas e independientes entre sí y promedia los espectros
  resultantes, reduciendo el ruido sin acortar la señal. De cada espectro se obtiene
  la potencia absoluta y relativa en cada banda, cocientes entre bandas
  y la pendiente con que decae el espectro hacia frecuencias altas, que
  cambia con la profundidad del sueño.

\item \textbf{Husos y ondas lentas con promediado de Woody.} Los husos se
  detectan filtrando el EEG central entre 11 y 16~Hz, calculando su
  envolvente y buscando tramos de 0,5 a 2~s que superen un umbral
  relativo al propio registro. Las ondas lentas se detectan
  en el EEG frontal con el criterio visto en clase:
  filtrado entre 0,2 y 4~Hz, pico negativo menor a $-60~\mu$V y amplitud
  pico a pico mayor a 75~$\mu$V. Para estimar la forma típica de las
  ondas lentas de cada paciente se las promedia, pero el ruido corre el
  pico detectado de cada una y el promedio resulta
  más ancho y chico que la onda real. El método de Woody
  \cite{woody} corrige esto de forma iterativa: desplaza cada onda
  hasta que se parezca lo más posible al promedio, recalcula el promedio
  con las ondas alineadas y repite hasta que los desplazamientos dejan de
  cambiar. Por época se obtienen la cantidad de husos y de ondas lentas,
  su amplitud y cuánto se parece cada onda a la forma típica.

\item \textbf{Filtrado espacial CSP.} Hasta acá cada canal se analiza
  por separado. CSP \cite{ramoser} busca una suma ponderada de los cuatro
  canales de EEG cuya potencia sea lo más grande posible en un estadio y
  lo más chica posible en otro. Se aplicará a los pares más confundidos
  (por ejemplo, N1 y REM).

\item \textbf{EOG, EMG y HRV.} Del EOG se toma su energía, que crece con
  los movimientos oculares de REM, y del EMG su amplitud, que cae con la
  atonía. Para la HRV se arma, con los latidos de la etapa 4, la serie de
  intervalos entre latidos, descartando los latidos ectópicos.
  En una ventana de 5~min centrada en cada época se mide la potencia en baja frecuencia
  (LF, 0,04--0,15~Hz) y en alta frecuencia (HF, 0,15--0,4~Hz): HF
  refleja la actividad parasimpática y aumenta en sueño profundo,
  mientras que el cociente LF/HF aumenta en REM \cite{taskforce}.

\item \textbf{Clasificación.} Usaremos análisis discriminante lineal
  (LDA), que asigna cada época al estadio cuyo perfil de descriptores es
  más parecido. Al ser lineal, el peso de cada descriptor en la decisión
  se puede leer directamente. El clasificador decide cada época por
  separado, pero el hipnograma es un continuo, y podría ser que asigne épocas
  que pertenezcan a distintos registros; por lo tanto, debemos elegir la secuencia de toda la noche que mejor combina lo que dice el
  clasificador con esas transiciones.

\item \textbf{Interfaz gráfica.} Aplicación en Streamlit para cargar un
  registro, ver cada canal en cada etapa y comparar el hipnograma
  estimado con el de los técnicos.
\end{enumerate}

\begin{table}[t]
\centering\small
\caption{Correspondencia con los contenidos de la consigna.}
\vspace{-0.6em}
\begin{tabularx}{\textwidth}{@{}p{0.34\textwidth}Lc@{}}
\toprule
Ítem de la consigna & Implementación & Etapa \\
\midrule
Método de adquisición & PSG SOMNOmedics, sistema 10-20 referido a mastoides, 256~Hz & Sec.~2.1 \\
Filtros convencionales y óptimos & Pasa-altos y notch FIR/IIR; Wiener y NLMS & 2, 3 \\
Detección de QRS y recorte en realizaciones & Pan-Tompkins; promedios por latido y por onda lenta & 4, 7 \\
Segmentación de señales & Segmentación adaptativa con modelo autorregresivo & 5 \\
Análisis de ritmos cerebrales & Potencia por banda, husos y ondas lentas & 6, 7 \\
HRV y periodograma de Lomb & Intervalos entre latidos, ectópicos, LF y HF & 9 \\
Interfaz gráfica & Aplicación en Streamlit & 11 \\
Técnicas no vistas en clase & Multitaper, Woody, CSP, modelo oculto de Markov & 6--8, 10 \\
\bottomrule
\end{tabularx}
\end{table}

% ============================================================
% 3. EVALUACIÓN
% ============================================================
\section{Evaluación}

Sobre los pacientes de prueba, la métrica principal es el $\kappa$ de
Cohen, comparable con el acuerdo entre técnicos (0,74) y con la red de
referencia (0,79). Se reportarán también la matriz de confusión y el F1 de cada estadio,
que combina cuántas épocas del estadio se detectan y cuántas de las
asignadas son correctas; es necesario porque N1 es poco frecuente y la
exactitud global puede ser alta aunque N1 se clasifique mal.

Para verificar las hipótesis: \textbf{H1}, se
compara la variación de la potencia por banda entre épocas consecutivas
de un mismo estadio con multitaper y con Welch; \textbf{H2}, se insertan
ondas lentas de forma conocida y con corrimientos al azar en EEG real de
vigilia y se mide qué tan bien recupera esa forma cada método;
\textbf{H3}, se compara el F1 de N1 y REM con y sin los descriptores
CSP; \textbf{H4}, se verifica que tras la remoción el promedio del EEG
alrededor de cada latido tienda a cero y se compara el $\kappa$ con y
sin estas etapas. Las comparaciones se hacen paciente por paciente con la
prueba de Wilcoxon, que no supone distribución normal.

% ============================================================
% BIBLIOGRAFÍA
% ============================================================
\section*{Bibliografía}
\footnotesize
\begin{thebibliography}{99}
\bibitem{alvarez} D.~Alvarez-Estevez, R.~M.~Rijsman, ``Inter-database validation of a deep learning approach for automatic sleep scoring,'' \emph{PLoS ONE}, 16(8):e0256111, 2021.
\bibitem{hmc} D.~Alvarez-Estevez, R.~M.~Rijsman, ``Haaglanden Medisch Centrum sleep staging database (v1.1),'' PhysioNet, 2022.
\bibitem{aasm} R.~B.~Berry et al., \emph{The AASM Manual for the Scoring of Sleep and Associated Events}, v2.4. AASM, 2017.
\bibitem{danker} H.~Danker-Hopfe et al., ``Interrater reliability for sleep scoring according to the R\&K and the new AASM standard,'' \emph{J. Sleep Res.}, 18(1):74--84, 2009.
\bibitem{trinder} J.~Trinder et al., ``Autonomic activity during human sleep as a function of time and sleep stage,'' \emph{J. Sleep Res.}, 10(4):253--264, 2001.
\bibitem{yasa} R.~Vallat, M.~P.~Walker, ``An open-source, high-performance tool for automated sleep staging,'' \emph{eLife}, 10:e70092, 2021.
\bibitem{thomson} D.~J.~Thomson, ``Spectrum estimation and harmonic analysis,'' \emph{Proc. IEEE}, 70(9):1055--1096, 1982.
\bibitem{prerau} M.~J.~Prerau et al., ``Sleep neurophysiological dynamics through the lens of multitaper spectral analysis,'' \emph{Physiology}, 32(1):60--92, 2017.
\bibitem{woody} C.~D.~Woody, ``Characterization of an adaptive filter for the analysis of variable latency neuroelectric signals,'' \emph{Med. Biol. Eng.}, 5:539--553, 1967.
\bibitem{ramoser} H.~Ramoser, J.~Müller-Gerking, G.~Pfurtscheller, ``Optimal spatial filtering of single trial EEG during imagined hand movement,'' \emph{IEEE Trans. Rehabil. Eng.}, 8(4):441--446, 2000.
\bibitem{widrow} B.~Widrow et al., ``Adaptive noise cancelling: principles and applications,'' \emph{Proc. IEEE}, 63(12):1692--1716, 1975.
\bibitem{he} P.~He, G.~Wilson, C.~Russell, ``Removal of ocular artifacts from electro-encephalogram by adaptive filtering,'' \emph{Med. Biol. Eng. Comput.}, 42:407--412, 2004.
\bibitem{pantompkins} J.~Pan, W.~J.~Tompkins, ``A real-time QRS detection algorithm,'' \emph{IEEE Trans. Biomed. Eng.}, 32(3):230--236, 1985.
\bibitem{donoghue} T.~Donoghue et al., ``Parameterizing neural power spectra into periodic and aperiodic components,'' \emph{Nat. Neurosci.}, 23:1655--1665, 2020.
\bibitem{molle} M.~Mölle, L.~Marshall, S.~Gais, J.~Born, ``Grouping of spindle activity during slow oscillations in human non-REM sleep,'' \emph{J. Neurosci.}, 22(24):10941--10947, 2002.
\bibitem{piantoni} G.~Piantoni et al., ``Individual differences in white matter diffusion affect sleep oscillations,'' \emph{J. Neurosci.}, 2013.
\bibitem{lomb} N.~R.~Lomb, ``Least-squares frequency analysis of unequally spaced data,'' \emph{Astrophys. Space Sci.}, 39:447--462, 1976.
\bibitem{taskforce} Task Force of the ESC and NASPE, ``Heart rate variability: standards of measurement, physiological interpretation and clinical use,'' \emph{Circulation}, 93(5):1043--1065, 1996.
\bibitem{rabiner} L.~R.~Rabiner, ``A tutorial on hidden Markov models and selected applications in speech recognition,'' \emph{Proc. IEEE}, 77(2):257--286, 1989.
\end{thebibliography}

\end{document}
